\section{How does sampling work?}
\label{sec:hws_how}

\begin{figure}
    \centering
    \resizebox{0.8\textwidth}{!}{%
    \begin{tikzpicture}[font=\small]
        \tikzstyle{normal}=[anchor=center, fill=white, draw]
        \tikzstyle{circ}=[normal, circle]
        \tikzstyle{dia}=[normal, diamond, aspect=1.5]
        \tikzstyle{linestyle}=[-{Latex[]}]

        \begin{scope}
            %%% background: user code vs sampling loop
            \fill[red, draw=red, fill opacity=0.2] (-3.9, 1.2) rectangle (3.4, -16.3);
            \node[red, anchor=east] at (3.1, -16) {user code};
            
            \fill[blue, draw=blue, fill opacity=0.2] (3.44, 1.2) rectangle (9.8, -16.3);
            \node[blue, anchor=west] at (3.7, -16) {asynchronous sampling loop};
        \end{scope}
    
        %%% main program
        \begin{scope}
            \node[normal] (construct) at (0, 0) (construct) {\code{system_hardware_sampler(@$\dots$@)}};
            \node[below = of construct, normal] (hws_start) {\code{start_sampling()}};
            \node[below =  of hws_start, circ] (thread_spawn) {spawn};
            \node[below right = .35cm and 1.65cm of thread_spawn, normal, dashed, fill=white!80!gray, anchor=east, rotate=90] (hws_pause) {\code{pause_sampling()}};
            \node[below left = 5.3cm and 1.5cm of hws_start, normal, dashed, fill=white!80!gray, rotate=90] (hws_event) {\code{add_event(@$\dots$@)}};
            \node[below = 8.0cm of thread_spawn, normal] (hws_stop) {\code{stop_sampling()}};
            \node[below = of hws_stop, circ] (join) {join};
            \node[above right = .35cm and .62cm of hws_stop, normal, dashed, fill=white!80!gray, anchor=west, rotate=90] (hws_resume) {\code{resume_sampling()}};
            \node[below = of join, normal] (hws_dump) {\code{dump_yaml(@$\dots$@)}};
            \node[below = of hws_dump, normal] (destruct) {\code{~system_hardware_sampler()}};
    
            \draw[linestyle, solid] (construct) -- (hws_start);
            \draw[linestyle, solid] (hws_start) -- (thread_spawn);
            \draw[linestyle, dashed] (thread_spawn.south) --  (hws_stop.north);
            
            \draw[linestyle, dashed] (thread_spawn.east) -- (hws_pause.east);
            \draw[linestyle, dashed] ([xshift=-0.1cm]hws_resume.east) -- ([xshift=-0.1cm]hws_pause.west);
            \draw[linestyle, dashed] ([xshift=0.1cm]hws_pause.west) -- ([xshift=0.1cm]hws_resume.east);

            \draw[linestyle, dashed] (construct) -| ([xshift=0.1cm]hws_event.east);
            \draw[linestyle, dashed] ([xshift=0.1cm]hws_event.west) |- (destruct);

            \coordinate[above left = 0.4cm and 0.5cm of hws_event.east] (hws_event_dummy_1);
            \coordinate[below left = 0.4cm and 0.5cm of hws_event.west] (hws_event_dummy_2);
            \draw[linestyle, dashed] ([xshift=-0.1cm]hws_event.east) |- (hws_event_dummy_1) -- (hws_event_dummy_2) -| ([xshift=-0.1cm]hws_event.west);

            \draw[linestyle, solid, thick] ([yshift=0.6cm]construct.north) -- (construct.north); 
            \draw[linestyle, solid, thick] (destruct.south) -- ([yshift=-0.6cm]destruct.south); 
            
            \draw[linestyle, dashed] (hws_resume.west) -- (hws_stop.east);
            \draw[linestyle, dashed] (hws_stop) -- (join);
            \draw[linestyle, solid] (join) -- (hws_dump);
            \draw[linestyle, solid] (hws_dump) -- (destruct);
        \end{scope}
    
        %%% sampling thread
        \begin{scope}
            \node[below right = 1.75cm and 5cm of thread_spawn, normal] (init_samples) {init samples};
            \node[below = 1.5cm of init_samples, dia] (decision) {\code{is_sampling}?};
            \node[right = 2cm of decision, circ, anchor=center] (wait) {wait};
            \node[below = 1.5cm of wait, normal, anchor=center] (sample) {sample};
    
            \draw[linestyle, solid] (init_samples) -- (decision);
            \draw[linestyle, solid] (decision.south) |- node[below, xshift=1.54cm] {sampling} (sample);
            \draw[linestyle, solid] ([yshift=0.1cm]decision.east) -- node[above] {paused} ([yshift=0.1cm]wait.west); 
            \draw[linestyle, solid] ([yshift=-0.1cm]wait.west) -- ([yshift=-0.1cm]decision.east);
            \draw[linestyle, solid] (sample) -- (wait);
        \end{scope}
    
        %%% paths between
        \draw[linestyle, dotted] (thread_spawn.east) -| (init_samples.north);
        \draw[linestyle, dotted] ([xshift=-.4cm, yshift=0.28cm]decision.south) |- node[left, xshift=-1.75cm, yshift=0.3cm] {stopped} (join.east);
        \draw[linestyle, dotted] (hws_pause.south) -- (decision.west);
        \draw[linestyle, dotted] (hws_resume.south) -- (decision.west);
        \draw[linestyle, dotted] (hws_stop.east) -| ([xshift=0cm]decision.west);
        
    \end{tikzpicture}
    }
    \caption{%
    An overview of the program flow of the hws hardware sampling library using the \cpp function names. 
    Constructing a hardware sampler initializes the respective vendor library environment. 
    A call to \code{start_sampling()} spawns a new \code{std::thread} of execution that contains the actual sampling loop. 
    This thread of execution is then notified and joined after a call to \code{stop_sampling()}. 
    Afterward, the gathered hardware samples can be saved to a YAML file. 
    During the application execution, hardware sampling can be paused and resumed using the respective functions. 
    Additionally, special points in the execution can be marked using custom events. 
    }
    \label{fig:program_flow}
\end{figure}

While the last section was concerned with \textit{what} the hws library can sample, this section explains \textit{how} it is implemented. 
If one is only interested in how to use the hws library, the next \autoref{sec:hws_example} will showcase an example usage of the hws library. 

The general internal flow of the hws library looks as follows. 
When the first hardware sampler for a specific hardware platform is constructed, the required environment from the respective vendor library is initialized. 
Using an atomic reference count also ensures that the environments are only initialized once, even when multiple hardware samplers for the same hardware type are used. 
The hardware sampling can then be started using the \code{start_sampling()} member function, which internally spawns a new \code{std::thread}. 
This newly created asynchronous thread is responsible for the actual sampling process. 
Each hardware sampler spawns its own thread, so it is advisable to use as few hardware samplers as possible to reduce overheads. 
The \code{system_hardware_sampler} also starts as many threads as it encapsulates other hardware samplers. 

For each hardware sampler, the sampling process starts with initializing all supported hardware samples: the fixed samples to their final values and the sampled samples to their first values. 
This step is also used to determine which hardware samples are actually supported by the current hardware platform. 
All theoretically supporters samples are queried with the respective vendor functions. 
If the function returns with an error code, the value will be set to an \code{std::nullopt}, indicating that it can be ignored in future sampling attempts during the sampling loop. 

Afterward, the actual sampling loop is started, adding a new value to every supported sampled sample in each iteration. 
At the end of one sampling iteration, the current thread sleeps for the amount of time specified by the sampling interval before it starts to sample new values again. 
Using the \code{pause_sampling()} and \code{resume_sampling()} functions, the sampling loop can be temporarily paused using an atomic flag as a signaling variable. 

If the hardware sampling should be completely stopped instead of only paused, the \code{stop_sampling()} member function must be called, which waits until the current sampling loop has finished and then joins on the previously spawned \code{std::thread}. 
Note that once a hardware sampler has been stopped, it can not be started again. 
Only after a hardware sampler has been stopped can all samples be saved to a YAML file using the \code{dump_yaml()} function. 
Note that it is still possible to retrieve the current hardware samples programmatically during the application's execution. 
At the end of the hardware sampler's lifetime, its destructor decreases the previously mentioned reference count. 
Only after the reference count reaches zero is the respective vendor-specific environment shut down, ensuring this happens only once. 

During the whole hardware sampling, the \code{add_event()} member function can be used to mark special points in the execution. 
This function accepts an event name, which is internally stored alongside the respective timestamp. 

\subsection{Implementation on GPUs}
\label{sec:hws_impl_gpus}

For implementing hws on the supported \glspl{gpu}, vendor-specific libraries are used. 
To sample information on NVIDIA \glspl{gpu} \gls{nvml}~\cite{nvidia_nvml_2024} is used, on \gls{amd} \glspl{gpu} the \gls{rocmsmi} library~\cite{amd_rocm_2024}, and on Intel \glspl{gpu} the Level Zero library~\cite{intel_level_2024}. 
While the \gls{nvml} and \gls{rocmsmi} \glspl{api} are relatively similar, Level Zero has a significantly different approach that can also be seen in \autoref{code:vendor_library_calls}. 
All library functions return a status code indicating whether the \gls{api} call was successful. 
If requested during CMake configuration, hws internally checks these status codes and raises an error if a function fails.  
Additionally, the \gls{nvml} and Level Zero specific device handles are wrapped in a custom wrapper class to prevent implementation details from leaking into the public headers. 

\subsection{Implementation on CPUs}
\label{sec:hws_impl_cpus}

On the \gls{cpu}, however, this is different since there is no simple library that can be called to gather all the necessary information. 
Therefore, three different command-line utilities are used if and only if they can be found during CMake configuration; otherwise, all related code is ignored, and the respective samples are not gathered. 
To retrieve general \gls{cpu} information, like the core count, socket count, or cache size, \texttt{lscpu}~\cite{qian_lscpu_nodate} is used, for memory-related information, i.e., available, used, and total main memory, the \texttt{free}~\cite{small_free_nodate} utility is used, and for everything else, like frequency or power related information, \texttt{turbostat}~\cite{brown_turbostat_nodate} is used. 
In all three cases, the \textit{subprocess.h} library~\cite{henning_subprocessh_2024} is used to create a new subprocess that, in turn, executes the respective utility, returning the standard output as a string, which is then post-processed, extracting the necessary information inside the \gls{cpu} hardware sampler.  

% lscpu
In case of the \texttt{lscpu} utility, the subprocess simply calls \texttt{lscpu}. 
Note that the cache sizes are only reported as strings instead of integers containing the actual values since the \texttt{lscpu} output is not standardized and varies greatly between \glspl{cpu}. 
%
% free
For the \texttt{free} utility, the command-line flag \texttt{-b} is added to ensure that the memory sizes are output in bytes. 

%turbostat
For \texttt{turbostat}, it is a bit more difficult. 
First, it should be noted that to get the full \texttt{turbostat} information, it has to be executed with root privileges. 
However, in newer \texttt{turbostat} versions, it is possible to get a subset of all possible information even if \texttt{turbostat} is not called with root privileges. 
The hws library handles this distinction during CMake configuration. 
If \texttt{turbostat} can be used with root privileges, the command \code{sudo turbostat -n 1 -i 0.001 -S -q} is executed; if root privileges are not granted, \code{sudo} is omitted. 
This command samples all available information immediately (\code{-i 0.001}) exactly once (\code{-n 1}) and outputs the information in a short summary format (\code{-S}), skipping potential header information (\code{-q}). 
An example \texttt{turbostat} output with the mentioned command is displayed in \autoref{code:turbostat_output}. 
\begin{listing}
    \begin{moderncode}{text}
Avg_MHz  Busy%  Bzy_MHz  TSC_MHz  IPC   IRQ  POLL  C1   C2   POLL%  C1%   C2%     CorWatt  PkgWatt
39       1.77   2212     4686     0.92  818  1     156  819  0.01   3.03  110.73  4.58     264.06
    \end{moderncode}
    \caption{%
    An example output of the turbostat command \code{sudo turbostat -n 1 -i 0.001 -S -q} with root privileges enabled on an \gls{amd} EPYC 9274F \gls{cpu}.
    }
    \label{code:turbostat_output}
\end{listing}
Note that the output of \texttt{turbostat} also depends on the used \gls{cpu}. 
For example, the output in \autoref{code:turbostat_output} was gathered on an \gls{amd} EPYC 9274F \gls{cpu} and consists of $14$ table entries. 
The same command executed on an Intel i5-10210U \gls{cpu}, the output consists of $52$ columns, mostly adding more idle state information but also temperature-related information missing for the \gls{amd} \gls{cpu}.

